% Part: turing-machines
% Chapter: machines-computations
% Section: variants

\documentclass[../../../include/open-logic-section]{subfiles}

\begin{document}

\olfileid{tur}{mac}{var}
\olsection{ٹیورنگ مشیناں دے وکھرے روپ}

ٹیورنگ مشیناں دی تعریف کرن دے کئی طریقے
ہو سکدے نیں؛ ساڈی تعریف اوہناں وچوں
صرف اک اے۔ کجھ گلاں وچ ساڈی تعریف
ہور تعریفاں توں زیادہ کھلی اے۔ اسی
نشاناں دا کوئی وی محدود مجموعہ پٹی
لئی منّ لیندے آں؛ اک زیادہ پابند تعریف
صرف دو نشان، $\TMstroke$
تے~$\TMblank$، منّ سکدی اے۔ اسی
مشین نوں پٹی اُتے نشان لکھن تے
نال ہی سر ہلاؤن دی اجازت دیندے آں؛
ہور تعریفاں وچ اک قدم اُتے
یا نشان لکھیا جا سکدا اے یا سر ہلایا جا سکدا اے۔
ساڈی تعریف وچ سر نوں ہلائے بنا
نشان لکھنا وی ممکن اے؛ کجھ ہور
تعریفاں $\TMstay$ والے “حکم” نوں
شامل نہیں کردیاں۔ دوجے پاسے کجھ
گلاں وچ ساڈی تعریف زیادہ پابند اے۔
اسی پٹی نوں صرف اک پاسے ول
لامتناہی منّیا اے؛ ہور تعریفاں وچ
پٹی کھبے تے سجے، دوہاں پاسیاں
ول لامتناہی ہو سکدی اے۔ بلکہ
وکھ وکھ پٹیاں دی کوئی وی تعداد،
یا خانیاں دی لامتناہی جالی وی
منّی جا سکدی اے۔ اسی ٹیورنگ
مشین دے حکماں نوں حالتی تبدیلی
والے فنکشن نال ظاہر کردے آں؛
ہور تعریفاں حالتی تبدیلی والا
تعلق ورتدیاں نیں، جیہدے وچ اکّو
صورتِ حال وچ مشین لئی اک توں
ودھ ممکن حکم ہو سکدے نیں۔

تعریف دی ایہ آخری ڈھیل خاص دلچسپ اے۔
ساڈی تعریف وچ جدوں مشین حالت~$q$
وچ نشان~$\sigma$ پڑھدی اے،
تاں $\delta(q, \sigma)$ نواں نشان،
نئی حالت تے سر دی اگلی جگہ
طے کردا اے۔ پر جے حکماں
دا مجموعہ موجودہ حالت-نشان
جوڑیاں $\tuple{q,
  \sigma}$ تے نویں حالت-نشان-سمت
تِکاں $\tuple{q', \sigma',
  D}$ دے وچکار تعلق ہووے،
تاں ٹیورنگ مشین دا عمل
اکّو طرح طے ہونا ضروری
نہیں۔ حکماں دے تعلق وچ
$\tuple{q,
  \sigma, q', \sigma', D}$ تے
$\tuple{q, \sigma, q'', \sigma'', D'}$
دوویں شامل ہو سکدے نیں۔
ایس صورت وچ مشین
\emph{غیر قطعی} (non-deterministic)
ٹیورنگ مشین ہوندی اے۔
حسابی پیچیدگی دے نظریے وچ
ایہناں دا اہم کردار اے۔

ٹیورنگ مشین کدوں رُکدی اے،
ایس بارے وی وکھرے رواج نیں:
ساڈے ہاں جدوں حالتی تبدیلی
والا فنکشن غیر تعریف شدہ ہووے
تاں مشین رُکدی اے؛ ہور تعریفاں
خاص مقرر رُکّن والی حالت
لازم کردیاں نیں۔ \olref[hal]{sec}
وچ اسی سمجھایا اے کہ مقرر
رُکّن والی حالت دی شرط
ٹیورنگ مشیناں دے حساب کرن
دے دائرے نوں نہیں گھٹاندی۔
ساڈیاں مشیناں دی پٹی صرف
اک پاسے ول لامتناہی اے،
ایس لئی کدی مشین کوئی
حکم ٹھیک طرح پورا نہیں
کر سکدی: جے سر سب توں
کھبے خانے نوں پڑھ رہیا
ہووے تے کھبے ول جانا
ہووے۔ ساڈی تعریف وچ سر
پٹی توں “ڈگّن” دی بجائے
اوہتھے ہی رہندا اے؛ پر
اسی ایہ وی تعریف کر سکدے
سن کہ ایس صورت وچ مشین
رُک جاوے۔ ایہ تعریف وی
ہم ارز اے: خانہ~$0$ توں
کھبے جان دی کوشش اُتے
رُکّن والی مشین دے عمل
دی نقل کرن لئی ہر اوہ
حالتی تبدیلی
$\delta(q,\TMendtape) =
\tuple{q',\sigma,\TMleft}$
حذف کر دیو۔ فیر
$\TMendtape$ اُتے کھبے
جان دی کوشش کرن دی
بجائے مشین رُک جاوے
گی۔\footnote{ایہ گل
\emph{پوری طرح} درست نہیں،
کیوں جے ساڈے قاعدیاں
وچ $\TMendtape$ نشان
خانہ~$0$ توں
علاوہ ہور خانیاں اُتے
وی لکھن
تے پڑھݨ توں کوئی
نہیں روکدا
(ویکھو \olref[una]{ex:mover})۔
ایس مشکل توں بچن لئی
ایس کم واسطے اک دوجا
نشان~$\TMendtape'$
شامل کیتا جا سکدا اے۔}

اعداد نوں ظاہر کرن دے
طریقے وی وکھرے ہو سکدے
نیں، ایس لئی ٹیورنگ مشین
دا حساب کیتا ہویا انپٹ-آؤٹ پٹ
فنکشن وی وکھرے روپ وچ
پیش ہو سکدا اے۔ اسی
اعداد دی اک نشانی نمائندگی
ورتدے آں؛ ثنائی نمائندگی
وی ورتی جا سکدی اے۔
ایس لئی $\TMblank$
تے~$\TMendtape$ توں
علاوہ دو نشان لوڑیں دے نیں۔

ہن اک دلچسپ حقیقت ویکھو:
ایہناں وکھرے روپاں وچوں
کسی روپ نال وی ایہ گل نہیں
بدلدی کہ کیہڑے فنکشن
ٹیورنگ مشین نال حساب ہو
سکدے نیں۔ \emph{جے کوئی
فنکشن اک تعریف مطابق
ٹیورنگ-حسابی اے، تاں اوہ
ایہناں ساریاں تعریفاں
مطابق ٹیورنگ-حسابی اے۔}

اسی ایس گل دی پوری
تصدیق دیاں تفصیلاں وچ
نہیں جاواں گے۔ بس اک
مثال ویکھو: دوہاں پاسیاں
ول لامتناہی پٹی یا کئی
پٹیاں دی اجازت نال
حسابی طاقت نہیں ودھدی۔
وجہ، موٹے طور اُتے،
ایہ اے کہ اکّو، اک
پاسے ول لامتناہی پٹی
والی ٹیورنگ مشین کئی
پٹیاں یا دوہاں پاسیاں
ول لامتناہی پٹی دی
نقل کر سکدی اے۔ مثال
دے طور اُتے، ودھ
حالتیں تے حکم ورت کے
کئی پٹیاں والی یا
دوہاں پاسیاں ول
لامتناہی پٹی والی
مشین دے پروگرام نوں
اک، اک پاسے ول
لامتناہی پٹی والی
مشین دے پروگرام وچ
“ترجمہ” کیتا جا
سکدا اے۔ نئی مشین
پٹی~$1$ دے خانیاں
(یا دوہاں پاسیاں
والی پٹی دے “مثبت”
خانیاں) لئی جُفت
خانے، تے پٹی~$2$
دے خانیاں (یا
“منفی” خانیاں)
لئی طاق خانے
ورت سکدی اے۔

\end{document}
